\section*{Erd\H{o}s Problem 1207 --- GPT 6 Astra Ultra partial attempt}


\subsection*{FORMAL RESTATEMENT}
For fixed dimension $d$, define $P_d(n)$ as the greatest integer $m$ such that every $n$-point subset of $\mathbb R^d$ contains an $m$-point subset without an isosceles triple. Estimate $P_d(n)$. In particular, does some absolute $c>0$ satisfy
\[P_2(n)<n^{1-c}\]
for every sufficiently large $n$?

\subsection*{CONTEXT AND SCOPE}
This replaces the earlier statement-only record with an elementary mathematical attempt. The arguments below establish only their stated partial results; no novelty or full solution is claimed.

The discussion identifies the one-dimensional case with avoidance of three-term arithmetic progressions. Thus isosceles triple includes three distinct collinear points when two of their mutual distances coincide; it is not restricted to nondegenerate triangles. The current official remarks credit Lee, Pohoata, and Zhu [LPZ26] with proving $P_2(n)<n^{1-c}$ for some $c>0$: this highlighted subquestion is resolved, while the general estimation problem remains open.

\subsection*{ATTACK PLAN AND WORK}
The collinear case admits a direct greedy lower bound. Identify a line with
the real numbers and let $S$ be an isosceles-triple-free subset of size $m$.
For each pair $a,b\in S$, a new point creates a three-term arithmetic progression
only at
\[x=(a+b)/2,\qquad x=2a-b,qquad x=2b-a.\tag{1}\]
These are exactly the collinear isosceles possibilities for three distinct points.
If $S$ is maximal in an $n$-point collinear set, (1) therefore gives
\[n\le m+3\binom m2,\qquad
m\ge\frac{1+\sqrt{1+24n}}6.\tag{2}\]
The bound follows by solving the resulting quadratic inequality, with integer
rounding understood. It proves a square-root lower bound for $P_1(n)$ and for
any higher-dimensional input set that happens to lie on a line.

This point-count argument cannot be copied unchanged in the plane. Equal
distances from a new point to two old ones constrain it to a whole perpendicular
bisector; equal distances from an old point constrain it to a circle. Either
curve can contain many points of the original set, so the three-forbidden-points
bound per pair is no longer valid. The official page already records a solution
of the highlighted planar upper-bound subquestion; this restricted lower-bound
argument makes no claim to reprove that result.

\subsection*{VERIFICATION AND REMAINING GAP}
Estimate the extremal function for unrestricted configurations in fixed
dimension, dealing with concentration on circles and bisectors. The proved
one-dimensional bound does not resolve the general estimation problem.

\subsection*{SOURCES}
https://www.erdosproblems.com/1207

https://www.erdosproblems.com/history/1207

\subsection*{FINAL}
\textbf{UNRESOLVED.} The full source problem remains outside the scope of the proved partial results.

\subsection*{COMPLETION ESTIMATE}
$5\%$. This is a conservative subjective estimate toward the whole problem, not a probability that the argument is correct or a claim that a fixed fraction of the problem has been solved.
